\chapter{给自己做一张资产负债表}
\chaptermeta{19}

\begin{ledetext}
第 5 章那三个词，汇丰在用，你们楼下便利店也在用。这一章读完要动手，不动手等于没读。
\end{ledetext}

\section{林悦不知道自己有多少钱}

林悦 32 岁，在杭州一家做企业 SaaS 的公司当产品经理。税前 2 万 6，五险一金扣完到手 1 万 9 出头。2021 年秋天她买了房，总价 386 万，首付 116 万，贷了 270 万，三十年。

过年饭桌上有亲戚问她现在有多少身家。她说：一套房，还欠着银行。

这个回答等于没说。

第 5 章那个式子是\textbf{资产 = 负债 + 净值}，全世界的会计都在用它，它不挑规模。套在林悦身上，需要的只是两列数字和一次减法。

唯一的技术难点在怎么填房子那一栏。买入价 386 万是历史，跟今天没有关系。\hlmark{房子按现在能卖掉的价钱记，房贷按银行系统里的余额记。}她翻了小区最近半年的成交记录，同户型同朝向，328 万落槌。就用 328 万。

\begin{bsheet}
  \bsside{资产 ASSETS}{%
    \bsrow{活期存款}{8,000}
    \bsrow{货币基金}{4.2 万}
    \bsrow{股票型基金}{11 万}
    \bsrow{公积金账户}{9.6 万}
    \bsrow{自住房（2026 年 8 月市价）}{328 万}
    \bsrowtotal{合计}{353.6 万}
  }
  \bsside{负债 + 净值}{%
    \bsrow{住房贷款余额}{247 万}
    \bsrow{信用卡未还}{1.3 万}
    \bsrow{车贷}{0}
    \bsrow{净值}{105.3 万}
    \bsrowtotal{合计}{353.6 万}
  }
\end{bsheet}

105.3 万。这是林悦真正拥有的东西，是把所有杠杆抽掉之后剩下的那一块。

接着是难受的一段算术。

\begin{egbox}{房子这一栏，五年下来是赚是赔}
林悦在这套房上投进去的本金是：首付 116 万，加上五年里还掉的本金 23 万（贷款从 270 万降到 247 万），合计 \textbf{139 万}。同期还掉的七十多万利息不算在内，那是买了五年的居住权。

今天这套房属于她的部分是：328 万市价减去 247 万余额，\textbf{81 万}。

139 万进去，81 万出来。这一栏五年少了 58 万。

她的净值总共 105.3 万，其中房子贡献 81 万。\textbf{她一直觉得房子是自己最成功的一笔投资，实际上是亏得最多的那一笔。}

\end{egbox}

\begin{warnbox}{小心：这个差额可以是负数}
林悦算是运气好的。换一个 2021 年在环沪、环郑、环沈那些地方接盘的人，跌四成不稀奇：386 万的房子如果今天只能卖 230 万，而贷款余额还有 247 万，资产负债表上这一栏就是 \textbf{−17 万}。

整张表的净值可能因此变成负的。人还是那个人，工作还是那份工作，只是表翻过去了。这个状态有个名字，叫\term{资不抵债}{negative equity}\index{092@资不抵债}。2008 年之后的美国有上千万套住房处在这个状态里，2026 年的中国有多少，没有公开统计。

不做这张表，你不会知道自己在不在里面。

\end{warnbox}

\begin{mythbox}{常见误解}
\mvfrow{误解}{"我有一套房，所以我有钱。"}
\mvfrow{事实}{房子记在资产端，房贷记在负债端，属于你的只有中间那个差。而且这个差有三个毛病：不能吃，不能应急，不能切一块下来卖。林悦手里能在 24 小时内变成钱的东西一共 5 万，她的表上却写着 353.6 万的资产。\textbf{净值是身家，流动性才是生活。}}
\end{mythbox}

\section{那张表是照片，这张才是心电图}

资产负债表拍的是某一天的静态。真正决定你三年后长什么样的，是每个月流过你账户的那道水。

结构简单到不需要软件：\textbf{收入 − 刚性支出 − 弹性支出 = 储蓄}。刚性支出是砍不掉的（月供、房租、保费、给父母的钱），弹性支出是失业第二个月就会自动消失的（旅游、外卖、健身卡、想买的那件外套）。

林悦一直觉得自己存不下钱。她把数字列出来之后愣了一下。

\begin{tablewrap}
\begin{xltabular}{\linewidth}{>{\raggedright\arraybackslash}p{0.20\linewidth}Xr}
\toprule
\bthead{项目} & \bthead{说明} & \bthead{每月} \\
\midrule
\textbf{工资到手} & 税后打到卡上的 & 19,000 \\
\textbf{公积金进账} & 单位 12\% + 个人 12\%，基数 2.6 万 & 6,240 \\
\textbf{房贷月供} & 刚性。剩 25 年，执行利率 3.35\% & −12,170 \\
\textbf{给父母} & 刚性 & −2,000 \\
\textbf{保险} & 刚性。一份 316 元/年的意外险，摊到月 & −26 \\
\textbf{弹性支出} & 吃饭通勤、健身卡、衣服、旅游摊销 & −5,800 \\
\textbf{结余} & 储蓄率 20.8\% & 5,244 \\
\bottomrule
\end{xltabular}
\end{tablewrap}

顺便说个真实细节：她 2021 年签的合同利率是 5.65\%，月供 15,585。经过 2023 年 9 月和 2024 年 10 月两轮存量房贷利率批量下调，现在执行 3.35\%，月供 12,170。这三千多块钱是政策给的，她自己什么都没做，甚至一开始没注意到。

更关键的是那个 6,240。公积金是她的收入，但她感觉不到，因为那笔钱直接进了另一个账户，还没经过她的手就被逐月提取去还了贷。\hlmark{感觉不到的收入，往往也是唯一存下来的收入。}

\begin{pullquote}{}
收益率不归你管，储蓄率归你管。它是这本书里唯一一个你能百分之百控制的变量。
\end{pullquote}

\begin{mythbox}{常见误解}
\mvfrow{误解}{"先想办法投资，光靠攒钱太慢了。"}
\mvfrow{事实}{算一遍就知道谁快。假设年化 5\%，每月存 2,500 块，存十年，终值 38.8 万。\\ 现在给你两个选项。\textbf{选项一：把收益率从 5\% 提到 8\%。}这件事极难，需要你连续十年跑赢绝大多数专业机构。做到了，终值 45.7 万，多出 \textbf{6.9 万}。\\ \textbf{选项二：把月存从 2,500 提到 5,200。}这件事只需要你重新排一遍支出，收益率还是那个 5\%。终值 80.7 万，多出 \textbf{41.9 万}。\\ 6.9 万对 41.9 万。前十年，储蓄率的权重压倒性地高于收益率。等你的本金攒到七位数，两者的地位才开始互换。}
\end{mythbox}

\section{三个数，自己量一下}

下面这几个门槛是经验值，不是定律，没有哪条监管规定写着 40\%。它们的作用像体检报告上的参考区间：超出去不代表你有病，代表该做个进一步检查。

\textbf{应急储备，按刚性支出算，不按总支出算。}林悦的刚性是月供 12,170 加赡养 2,000 加保费 26，再加上吃住通勤里怎么也省不掉的 2,800，大约 17,000。三个月 5.1 万，六个月 10.2 万。她手上活期 8,000 加货基 4.2 万，正好 5 万，卡在下限。

为什么用刚性支出？因为失业的第二个月你不会去旅游，不会续健身卡，弹性支出的定义就是它能归零。用总支出算会把这个数虚高一截，看着安心而已。

放的地方也有讲究：得是当天能到账的。货基快赎单日有一万的额度上限，七天通知存款、国债逆回购到期都行。不能放在 R2 以上的理财里，更不能放在需要 T+7 才回到卡上的产品里。急用钱的时候差的往往不是钱，是七天。

林悦另外那 11 万股票基金呢？那不是安全垫。那是一笔在你最需要用钱的时候，大概率也正好在跌的东西。经济不好的时候公司裁员和市场下跌是同一件事的两个面。

\textbf{月供收入比，看你怎么定分母。}林悦 12,170 除以到手的 19,000，是 64\%，吓人；把公积金算进收入，12,170 除以 25,240，是 48\%。同一个人，同一笔月供，两个结论差一倍。

这个指标最容易被自己骗过去。经验上超过 40\% 会开始难受，超过 50\% 基本没有抗风险能力，一次降薪或者三个月的空窗期就能击穿。林悦在 48\%，她自己觉得"还行"，因为公积金那部分她感觉不到。感觉不到不等于不存在。

\textbf{净值增长率，比收入增长率诚实。}今年涨薪 12\% 是件好事，但它不告诉你任何关于财富的信息，你完全可能同时把支出提高了 15\%。净值从年初的 92 万变成年末的 105 万，涨了 14\%，这个数才说明问题。它还有个好处：它会诚实地把房价跌的那部分写进去，不给你留面子。

\section{你最大的一笔资产，任何对账单上都没有}

林悦到手一年 23 万上下。她 32 岁，就算干到 60 岁，还有 28 年可以卖。这笔未来的收入有个名字，叫\term{人力资本}{human capital}\index{044@人力资本}。

\begin{egbox}{给自己估个价}
用第 6 章那套折现的做法：未来 28 年，每年 23 万，按 4\% 折回今天。

年金现值系数 =（1 − 1.04 的 −28 次方）÷ 0.04 = \textbf{16.66}。

23 万 × 16.66 = \textbf{383 万}。

折现率换成 3\%，这个数变成 432 万；换成 5\%，变成 343 万。区间不小，但结论不受影响：\textbf{383 万比她那套 328 万的房子还大，是她 105 万净值的三倍多。}

她全部资产里最大的一项，是她自己。而这一项没有对账单，没有净值曲线，也没人每天提醒她它今天跌了多少。

\end{egbox}

\begin{pullquote}{}
对绝大多数三十来岁的人来说，最大的一笔资产是自己，而且是唯一一笔没上保险的资产。
\end{pullquote}

把人力资本当成一只债券来看待，很多事情会立刻清楚。这只债券每月付息 1 万 9，到 60 岁那天本金归零，中间存在违约的可能。比喻可以继续推下去：面对一只有违约风险的债券，你会做两件事，一是给它买保险，二是不要在组合里再堆一大堆和它高度相关的东西。

第一件事推出的结论是：\textbf{年轻人真正的最大风险从来不是股票跌了。}林悦的基金全亏光是 11 万；她要是因为一场病三年不能工作，损失是三年 69 万，加上回来之后议价能力的下降。量级差六倍以上。可她花在挑基金上的时间，是花在体检和保单上的几十倍。这个精力分配是反的。

第二件事推出的结论更有意思：\textbf{你的人力资本长得像债券，还是像股票？}

\begin{tablewrap}
\begin{xltabular}{\linewidth}{>{\raggedright\arraybackslash}p{0.20\linewidth}XX}
\toprule
\bthead{职业} & \bthead{现金流的样子} & \bthead{和股市的相关性} \\
\midrule
\textbf{公务员、电网、三甲医院医生} & 稳定、可预测、涨得慢 & 接近零，像 AAA 级长期债券 \\
\textbf{企业 SaaS 产品经理（林悦）} & 看客户预算，客户预算看资本开支 & 中高，偏科技板块 \\
\textbf{券商营业部、投行、量化} & 底薪低，奖金和市场同呼吸 & 极高，像一只高贝塔股票 \\
\bottomrule
\end{xltabular}
\end{tablewrap}

在券商工作的人重仓券商股，等于把工资、奖金、股票账户、期权全押在同一个变量上。市场好的时候四个数一起涨，人会真心觉得自己有天赋；市场转向的时候，部门缩编、奖金归零、持仓腰斩会在同一个季度里发生。2015 年下半年和 2018 年，这种事在陆家嘴和金融街成批出现过。

林悦的客户预算和 AI 资本开支高度绑定。她的人力资本偏股票，还偏科技股。那 11 万基金里如果有 7 万是 AI 主题，她是在同一件事上加倍下注。

所以我给一个判断，而且认为它几乎没有例外：\textbf{大多数人的组合里最该被减配的，是自己公司的股票和自己行业的主题基金。}这跟看不看好那个行业无关，跟你已经在里面押了多少无关，跟你已经押了 383 万有关。

\section{四个盒子，按日期分，不按胆量分}

钱不该混在一个池子里。分开不是为了仪式感，是因为不同的钱有不同的到期日。

\begin{tablewrap}
\begin{xltabular}{\linewidth}{>{\raggedright\arraybackslash}p{0.20\linewidth}XXX}
\toprule
\bthead{账户} & \bthead{装什么} & \bthead{期限} & \bthead{放什么} \\
\midrule
\textbf{日常} & 一个月的开销 & 30 天内 & 活期、余额宝 \\
\textbf{安全垫} & 3–6 个月刚性支出 & 随时可能用 & 货基、七天通知存款、逆回购 \\
\textbf{目标} & 三年内确定要花的：首付、生育、读书、换车 & 1–3 年 & 定存、短债、同业存单指数 \\
\textbf{长期} & 十年不打算动的 & 10 年以上 & 权益类为主 \\
\bottomrule
\end{xltabular}
\end{tablewrap}

\begin{pullquote}{}
期限决定了你能承受多少波动，那张风险测评问卷决定不了。
\end{pullquote}

推一步就明白了。一个在问卷上被判定为"积极型"、但两年后要付 80 万首付的人，那 80 万碰都不该碰权益，因为两年不够一次熊市走完。一个被判定为"保守型"、但那笔钱三十年不会动的人，全放定存反而是在温水里慢慢亏掉购买力。同一份问卷，同一个人，两笔钱，答案完全相反。

\section{保险不是理财，保险是防破产}

这一节的全部内容就是一句话：\textbf{保险买的是"倒霉的时候不至于破产"，不是买收益。}

排序的原则也简单，先堵最可能发生并且最贵的那个窟窿。医疗险一年几百块，管大额住院，发生概率最高，放第一位。意外险一年一两百，杠杆最高，放第二位。重疾险赔一笔现金，用来填生病期间没有收入的那段，年缴几千，放第三位。定期寿险排最后，但它对某些人最要命。

\term{定期寿险}{term life insurance}\index{011@定期寿险}只在被保人身故时赔付，赔给还活着的人。它在中国被严重低估，因为买它的人自己一分钱好处都得不到。

林悦要是明年结婚生孩子，房贷还剩两百四十来万。她出事，这 247 万不会消失，会落在伴侣和孩子头上。\textbf{有房贷、有孩子的人，定期寿险保额至少要覆盖房贷余额。}她查了几家，32 岁女性、300 万保额、保到 60 岁、缴 20 年，年缴在 1,900 到 2,600 之间。一年两千块换掉一个 300 万的窟窿，是她整张表上性价比最高的一笔支出。

\begin{warnbox}{用保险做投资之前，先算清楚三件事}
\textbf{一，你交进去的钱有多少真的进了账户。}初始费用、保障成本、账户管理费会先拿走一部分，前几年拿走的比例通常最高。

\textbf{二，写进合同的保证利率是多少。}演示利率和保证利率是两个词，差一个字，差很多年。销售讲的通常是前者，赔付时按的是后者。

\textbf{三，什么时候能取，取的时候扣多少。}很多产品前五到十年退保拿回来的钱少于你交进去的。

这三个数全在合同里，全都不在话术里。如果有人花二十分钟讲收益、一分钟带过费用，你至少知道了他的重点在哪儿。

\end{warnbox}

\section{2026 年的算术题：要不要提前还贷}

\begin{statrow}{3}
  \statitem{3.5}{\%}{2026 年 8 月 20 日，5 年期以上 LPR}
  \statitem{1.4}{\%}{7 天期逆回购政策利率}
  \statitem{1.41}{\%}{商业银行净息差（2026 年二季度末）}
\end{statrow}

净息差 1.41\% 是历史低位。它的含义很直白：银行赚的存贷利差已经薄得像纸。利差压到这个程度，它就没有空间给你高存款利率了。所以 2026 年那种"我想找个稳当的地方拿 4\%"的愿望，基本无解。凡是在这个环境里同时承诺 4\% 以上和"保本"的东西，第 23 章会告诉你它的零件是怎么装的。

在这样的利率环境里，提前还贷成了很多人的常规问题。这里不给结论，给一套算法。

\begin{flowlist}
  \flowitem{01}{先查清楚你的房贷现在到底是多少利率}{不是合同上签的那个数。2023 年 9 月和 2024 年 10 月两轮存量利率调整之后，很多人的实际执行利率已经落到 3.1\%–3.5\% 区间，自己却还在按五年前的印象做判断。打开手机银行查重定价后的执行利率，五分钟的事。}
  \flowitem{02}{再找出你的"确定性收益率"}{不是你希望赚到的，是确定拿得到的：三年期定存、货基七日年化、国债逆回购。2026 年 8 月这一档大致在 1.4\% 到 2\% 之间。别把股票的期望收益放进这一格，期望不是确定。}
  \flowitem{03}{两个数相减}{3.35\% 减 1.8\%，差 1.55 个百分点。这才是提前还贷每年真正给你省下的东西，不是 3.35\%。很多人在这一步把收益算成了两倍。}
  \flowitem{04}{算完先别按确认}{还完之后你的安全垫还剩几个月？房子不能切一块下来救急，还进去的钱想拿回来只能重新申请贷款，而你需要它的那个时刻，往往正好是银行不愿意贷给你的时刻。\textbf{流动性只在你不需要它的时候，看起来像是多余的。}}
  \flowitem{05}{最后一格叫心理成本，没有公式}{月供从 12,170 降到 9,000，晚上睡得踏实一点，这件事值多少钱。有人填 0，有人填得比那 1.55 个百分点还大。两种人都没填错。}
\end{flowlist}

你的利率、你的备用金厚度、你工作的稳定性，我一样都不知道。\hlwarm{这一节不构成投资建议，也不构成还贷建议。}任何人在不知道这三件事的前提下告诉你"该不该提前还贷"，他给你的不是建议，是内容。

\bookbreak

最后说回动手这件事。这张表第一次做，会花掉你一个晚上，因为你得去翻公积金余额、查贷款剩余本金、找小区最近的成交价。第二次做只要二十分钟。

林悦做完之后改了三件事：把货基加到 8 万，买了那份定期寿险，把 AI 主题的仓位砍掉一半换成宽基。她的收益率没有因此提高，她只是不再需要靠猜测来回答"我到底有多少钱"。

\begin{takeaway}{本章一句话}
净值是算出来的，不是感觉出来的。而这张表上最大的一项，是你自己。

\begin{itemize}
  \item 房子按市价记，房贷按余额记，中间那个差才是你的，它可以是负数。
  \item 收益率不归你管，储蓄率归你管。前十年，后者的权重压倒性地大。
  \item 应急储备按刚性支出算 3–6 个月，放在当天能到账的地方。
  \item 32 岁的 383 万人力资本，比 105 万净值更需要被保护，也更容易被重复押注。
\end{itemize}

\end{takeaway}

\begin{bridgebox}
\textbf{下一章}表做完了，下一个问题绕不过去：那些数字要变大，钱到底得从哪儿来？
\end{bridgebox}

